\documentclass{silkmath-fixture}
\usepackage{amsmath,amssymb,booktabs,xcolor}
% 上面的 sty/cls 已由 \documentclass 加载：\eps \R \A \norm \Ocal \Tr 颜色、eqmath、question/solution。

\title{Silk Math 系统预览样例}
\begin{document}

\section{分隔符：\$ \$、\$\$ \$\$、\verb|\(| \verb|\)|、\verb|\[| \verb|\]|}

行内 $E=mc^2$ 与带宏的 $u\in\R^d$、$\norm{\A x}\le\eps$。

把光标放进行尾公式，预览应从公式下方往右，而不是行首：前面垫一些文字 $\norm{\A u}$。

圆括号行内 \(\alpha+\beta=\gamma\)，\(\pd{u}{t}=\Delta u\)。

独立美元段：
$$
\int_\Omega \abs{\nabla u}^2\,\mathrm{d}x = 0
$$

方括号独立段：
\[
\argmin_{x\in\R^n} \norm{\A x-b}_2^2
\]

行内公式跨行（中间不能有空行）：
$\frac{1}{2\pi}
\int_0^{2\pi} e^{in\theta}\,\mathrm{d}\theta$

\section{分隔符里再套顶层环境}

$$
\begin{equation}
  u_t + \operatorname{div}(u\A) = 0
\end{equation}
$$

\[
\begin{align}
  a &= b+c \\
  d &= e+f
\end{align}
\]

\section{公式环境}

\begin{equation}
  \label{eq:energy}
  E(u)=\int_\Omega \abs{\nabla u}^2\,\mathrm{d}x
\end{equation}

\begin{equation*}
  \mathrm{e}^{i\pi}+1=0
\end{equation*}

\begin{align}
  x &= y+z \label{eq:a} \\
  u &= \Tr(\A^\top\A)
\end{align}

\begin{align*}
  f(x) &= x^2 \\
  g(x) &= \sin x
\end{align*}

\begin{alignat}{2}
  x &= 1 &\qquad y &= 2 \\
  u &= 3 &\qquad v &= 4
\end{alignat}

\begin{gather}
  a+b=c \\
  d+e=f
\end{gather}

\begin{gather*}
  1+1=2 \\
  2+2=4
\end{gather*}

\begin{multline}
  a+b+c+d+e+f \\
  +g+h+i+j
\end{multline}

\begin{equation}
\begin{split}
  (a+b)^2
  &= a^2+2ab+b^2
\end{split}
\end{equation}

\begin{equation}
\begin{aligned}
  u &= \A v \\
  w &= \B u
\end{aligned}
\end{equation}

\section{类文件里的数学环境 vs 文本环境}

\begin{eqmath}
  \clsOnly{X} = \Ocal(h^2)
\end{eqmath}

\begin{question}
  证明 $u\equiv-1$。里面的行内公式和
  \[
    \Delta u = 0 \quad\text{in }\Omega
  \]
  都应单独预览，不能把整段 question 当成一条公式。
\end{question}

\begin{solution}
  由能量 $E(u)=0$ 得 $\nabla u=0$，因此 $u$ 为常数。
\end{solution}

\section{矩阵、array、cases}

\[
\begin{matrix} 1 & 2 \\ 3 & 4 \end{matrix}
\quad
\begin{pmatrix} a & b \\ c & d \end{pmatrix}
\quad
\begin{bmatrix} 1 & 0 \\ 0 & 1 \end{bmatrix}
\quad
\begin{Bmatrix} x \\ y \end{Bmatrix}
\quad
\begin{vmatrix} a & b \\ c & d \end{vmatrix}
\quad
\begin{Vmatrix} u \\ v \end{Vmatrix}
\]

\[
\begin{smallmatrix} 1&2\\3&4 \end{smallmatrix}
\qquad
\begin{array}{c|cc}
  a & b & c \\
  \hline
  d & e & f
\end{array}
\]

\[
  f(x)=\begin{cases}
    x^2 & x\ge 0 \\
    -x  & x<0
  \end{cases}
\]

\section{underbrace / overbrace / 根号 / 划掉}

\[
  \underbrace{a+b+c}_{\text{sum}} + \overbrace{xy}^{\text{prod}}
\]

\[
  \underbrace{\A u}_{=0 \text{ by the PDE}}
\]

\[
  \overline{z+w}\quad \underline{abc}\quad \sqrt{x^2+1}\quad \sqrt[3]{8}
\]

\[
  \hat a \quad \bar b \quad \vec{v} \quad \dot u \quad \ddot u
  \quad \tilde f \quad \widehat{ABC} \quad \widetilde{xyz}
\]

\[
  \cancel{x} \quad \bcancel{y} \quad \xcancel{z} \quad \cancelto{0}{h}
\]

\section{分数、上下标、binom、substack}

\[
  \frac{1}{2}+\dfrac{a}{b}+\tfrac{1}{n}+\frac12+\sqrt{\frac{p}{q}}
\]

\[
  x^{2}_{i,j} \quad a^{b^{c}} \quad \binom{n}{k}
  \quad \sum_{\substack{i=1\\ i\neq j}}^{n} a_{ij}
\]

\section{大小定界符}

\[
  \left( \frac{a}{b} \right)
  \quad \bigl[ x \bigr]
  \quad \Bigl\{ y \Bigr\}
  \quad \biggl\langle u \biggr\rangle
  \quad \left\{ x \;\middle|\; x>0 \right\}
\]

\section{文本、字体、颜色}

\[
  \text{energy } E(u)
  \quad \mathrm{d}x
  \quad \mathbf{x}
  \quad \mathbb{R}
  \quad \mathcal{L}
  \quad \mathscr{H}
  \quad \mathfrak{g}
  \quad \boldsymbol{\beta}
  \quad \bm{v}
\]

\[
  \textcolor{Accent}{u}
  +\textcolor{Ok}{v}
  +\textcolor{Warn}{w}
  +\textcolor{ClsAccent}{\clsOnly{Y}}
\]

\[
  \text{中文标签：} \quad u_{\text{内点}}
\]

\section{自定义宏（来自 .sty / .cls）}

\[
  \ip{\A u}{v} + \norm{\Ocal} + \pd{u}{x} + \diag(\B) + \Tr(\A)
\]

正文里再声明一次，当前文件宏也应生效：
\newcommand{\fileMacro}[1]{\mathrm{#1}}
\[
  \fileMacro{id} + \A
\]

\section{LaTeX 表格 tabular / longtable / booktabs}

\begin{tabular}{l|cr}
  左 & 中 & 右 \\
  $a$ & $b$ & $c$
\end{tabular}

\begin{tabular}{|c|c|c|}
  \hline
  A & B & C \\
  \hline
  1 & 2 & 3 \\
  \hline
\end{tabular}

\begin{tabular}{cc|c}
  \toprule
  方法 & $L^2$ 误差 & 阶 \\
  \midrule
  Crank--Nicolson & $3.1\times 10^{-4}$ & $2.00$ \\
  Backward Euler & $1.7\times 10^{-2}$ & $1.01$ \\
  \bottomrule
\end{tabular}

\begin{tabular}{cc}
  \multicolumn{2}{c}{合并单元格} \\
  $x$ & $y$
\end{tabular}

\begin{tabular}{*{3}{c}}
  a & b & c \\
  d & e & f
\end{tabular}

\begin{tabular}{p{3cm}c}
  自动换行的段落列 & $z$
\end{tabular}

\begin{longtable}{cc}
  头 & 值 \\
  $\alpha$ & 1 \\
  $\beta$ & 2
\end{longtable}

\section{长 aligned（滚动时预览应跟着可见行）}

\begin{align}
  a_1 &= 1 \\
  a_2 &= 2 \\
  a_3 &= 3 \\
  a_4 &= 4 \\
  a_5 &= 5 \\
  a_6 &= 6 \\
  a_7 &= 7 \\
  a_8 &= 8 \\
  a_9 &= 9 \\
  a_{10} &= 10
\end{align}

\section{光标容易踩坑的位置}

% 无花括号的分数参数、文本里的竖线、underbrace 参数缝
$\frac12+\frac1{2}$
$\text{for } x\in\R$
$\underbrace{abc}_{=0}$
$\textcolor{Accent}{u}$

\section{注释里的美元不应成公式}

% $this is a comment$
正文 $x$ 才是公式。

\section{未定义命令（应红字，其余仍渲染）}

\[
  \thisCommandDoesNotExist + x^2
\]

\section{TikZ（应提示需要完整 TeX 引擎，而不是空白）}

\begin{tikzpicture}
  \draw (0,0) -- (1,1);
\end{tikzpicture}

\end{document}
